% Prose for this counterexample.  Every region below is delimited by an
% environment that tools/build.py extracts; machine facts (ids, uid, dates,
% urls, enums) live in case.json.  See counterexamples/AGENTS.md.

\cxtitle{A mixed-norm Cauchy--Schwarz question of Sah, Sawhney, Stoner and Zhao}

% The question reached S.~Sra from Y.~Zhao in May 2018, months before the
% preprint (arXiv 1809.09462, September 2018), so the paper is where it was
% later stated in print rather than where it was first put.  Both readings below
% trace to that communication, which is why this sits in the case-level block
% and not in a per-result override.
\begin{cxcredits}
\posedby{A.~Sah, M.~Sawhney, D.~Stoner and Y.~Zhao \citeyearpar{SahSawhneyStonerZhao2020}, and put to S.~Sra by Y.~Zhao in private communication, May 2018}
\foundby{GPT-5.6 Sol (Pro)}{2026-07-13}
\formalizedby{S.~Sra}
\auditedby{S.~Sra}
\contributedby{S.~Sra}
\end{cxcredits}

% The two statements were refuted in different ways, so each overrides the
% shared block's \foundby: the general case is a proof whose explicit witness
% was generated by Opus 5, the $B=A$ case a witness produced by GPT-5.6 Sol.
\begin{cxcredits}[mixed-norm-general-s]
\foundby{Opus 5}{2026-08}
\end{cxcredits}

\begin{cxcredits}[rank-two-mixed-norm]
\foundby{GPT-5.6 Sol (Pro)}{2026-07-13}
\end{cxcredits}

\begin{cxcontext}
For a matrix $X$ the \emph{mixed norm} used below is
\[
\|X\|_{p,q}=\left(\sum_j\left(\sum_i |x_{ij}|^p\right)^{q/p}\right)^{1/q}.
\]
\citet{SahSawhneyStonerZhao2020} index the two sums the other way round, so their $\|X\|_{L_{p,q}}$ is $\|X^{\mathsf T}\|_{p,q}$ in the notation above; the distinction is immaterial for both statements below, since one is applied to a symmetric matrix and the other at $p=q$, where the mixed norm is just the entrywise $\ell_p$ norm.  A matrix is \emph{completely positive} when it is $Y^{\mathsf T}Y$ for some entrywise nonnegative $Y$; equivalently, it is the Gram matrix of finitely many vectors in the nonnegative orthant.  Their Lemma~3.1 --- proved by convex duality, an argument the present author supplied --- settles the case $p=1$, and a Schur-product argument settles every positive integer $p$; what remained open was every non-integer $p$.
\end{cxcontext}

% ---------------- result: mixed-norm-general-s ----------------

\begin{cxsource}{mixed-norm-general-s}
A.~Sah, M.~Sawhney, D.~Stoner and Y.~Zhao, the open question stated immediately after inequality~(3.1) in \S3
\end{cxsource}

\begin{cxstatement}{mixed-norm-general-s}
We do not know if the inequality can be extended to $\|A^{\mathsf T}B\|_{L_{p,q}}^2\le\|A^{\mathsf T}A\|_{L_{q,q}}\|B^{\mathsf T}B\|_{L_{p,p}}$ for all reals $1\le p\le q$ (this is true for positive integer $p$ by a tensor-power argument).
\end{cxstatement}

\begin{cxsummary}{mixed-norm-general-s}
Extension of the mixed-norm Cauchy--Schwarz inequality to all real $1\le p\le q$ \citep{SahSawhneyStonerZhao2020}
\end{cxsummary}

\begin{cxcertificate}{mixed-norm-general-s}
Nonnegative integer $A,B\in\mathbb Z_{\ge0}^{3\times2}$ at $p=q=3/2$; the deficit is an integer combination of integer square roots, negative by exact algebraic evaluation.
\end{cxcertificate}

% ---------------- result: rank-two-mixed-norm ----------------

\begin{cxsource}{rank-two-mixed-norm}
A.~Sah, M.~Sawhney, D.~Stoner and Y.~Zhao, the same open question at $B=A$; put to S.~Sra by Y.~Zhao in private communication, May 2018
\end{cxsource}

\begin{cxstatement}{rank-two-mixed-norm}
Taking $B=A$ in the question above and writing $Z=A^{\mathsf T}A$ for the resulting completely positive matrix: for all reals $1\le s\le q$ and every completely positive $Z$,
\[\|Z\|_{s,q}^2\le\|Z\|_{s,s}\|Z\|_{q,q}.\]
\end{cxstatement}

\begin{cxsummary}{rank-two-mixed-norm}
Rank-two completely-positive mixed-norm interpolation at $(s,q)=(6/5,6)$
\end{cxsummary}

\begin{cxcertificate}{rank-two-mixed-norm}
Explicit 21-vector Gram matrix and interval-separated ratio $>1$.
\end{cxcertificate}

\begin{cxrefutation}
The question fails in both of the readings above, by witnesses of quite different character.  That it has to fail at every non-integer exponent is a consequence of the Hadamard-power theory of \citet{FitzGeraldHorn1977} and \citet[Theorem~4.3]{HornMathias1990}; that argument is recorded as Remark~\ref{rmk:fh} below, and it is not new here.  What Theorem~\ref{thm:mixed-general} adds is a minimal explicit witness for it --- small enough to check by hand, and exact enough to need no enclosure.  The second witness is larger and interval-certified, but it refutes the harder $B=A$ specialization, which that mechanism cannot reach.

\subsubsection*{The general question}

\begin{theorem}[\statusfalse: extension to non-integer exponents]\label{thm:mixed-general}
Let $p=q=3/2$ and put
\[
A=\begin{bmatrix}2&4\\5&0\\1&0\end{bmatrix},\qquad
B=\begin{bmatrix}0&5\\4&2\\1&0\end{bmatrix}\ \in\ \mathbb Z_{\ge0}^{3\times2}.
\]
Then $\|A^{\mathsf T}B\|_{p,q}^2>\|A^{\mathsf T}A\|_{q,q}\|B^{\mathsf T}B\|_{p,p}$.
\end{theorem}
\begin{proof}
Both matrices are entrywise nonnegative, so they satisfy the hypothesis of the question as posed, not merely the weaker variant in which only the three products are required to be nonnegative.  Those products are
\[
A^{\mathsf T}A=\begin{bmatrix}30&8\\8&16\end{bmatrix},\qquad
B^{\mathsf T}B=\begin{bmatrix}17&8\\8&29\end{bmatrix},\qquad
A^{\mathsf T}B=\begin{bmatrix}21&20\\0&20\end{bmatrix}.
\]
At $p=q=3/2$ the mixed norm is the entrywise $\ell_{3/2}$ norm, so $\|M\|_{3/2,3/2}^{3/2}=\sum_{ij}|m_{ij}|^{3/2}$ and the deficit
\[
\Delta:=\|A^{\mathsf T}A\|_{3/2,3/2}^{3/2}+\|B^{\mathsf T}B\|_{3/2,3/2}^{3/2}-2\|A^{\mathsf T}B\|_{3/2,3/2}^{3/2}
\]
is an integer combination of square roots of integers:
\[
\Delta=64+64\sqrt2+30\sqrt{30}+17\sqrt{17}+29\sqrt{29}-42\sqrt{21}-160\sqrt5 .
\]
Deciding the sign of $\Delta$ is exact algebra over a real field, with no rounding anywhere:
\[
\Delta=-5.150045302282868288532184137690370433502833719528\ldots<0 .
\]
Hence $2\|A^{\mathsf T}B\|_{3/2,3/2}^{3/2}>\|A^{\mathsf T}A\|_{3/2,3/2}^{3/2}+\|B^{\mathsf T}B\|_{3/2,3/2}^{3/2}$, and the arithmetic--geometric mean inequality bounds the right side below by $2\bigl(\|A^{\mathsf T}A\|_{3/2,3/2}\|B^{\mathsf T}B\|_{3/2,3/2}\bigr)^{3/4}$.  Cancelling and raising to the power $2/3$,
\[
\frac{\|A^{\mathsf T}B\|_{3/2,3/2}}{\|A^{\mathsf T}A\|_{3/2,3/2}^{1/2}\|B^{\mathsf T}B\|_{3/2,3/2}^{1/2}}=1.00629360763001504125\ldots>1 . \qedhere
\]
\end{proof}

\begin{remark}[Why non-integer exponents were always going to fail]\label{rmk:fh}
The witness is not an accident of $3/2$.  Let $s\notin\mathbb N$ with $s<n-2$.  By the sharpness half of the FitzGerald--Horn theorem \citep{FitzGeraldHorn1977} there is a positive semidefinite $Z\in\R_{\ge0}^{n\times n}$ whose Hadamard power $Z^{\odot s}$ is not positive semidefinite, so $\sigma^{\mathsf T}Z^{\odot s}\sigma<0$ for some $\sigma$.  When $\sigma$ can be taken in $\{\pm1\}^n$ --- as it can here, because rescaling $Z$ by a positive diagonal is free --- splitting the index set along the sign of $\sigma$ writes $Z$ as the block Gram matrix of two nonnegative blocks $A,B$, and the displayed quadratic form is exactly $-\Delta$.  This is the mechanism of \citet[Theorem~4.3]{HornMathias1990}, and it produces a counterexample for \emph{every} non-integer $s$, which is why the positive results stop precisely at the integers.
\end{remark}

\subsubsection*{The specialization $B=A$}

\begin{theorem}[\statusfalse; computer-assisted exact construction]\label{thm:mixed}
There is an explicit $21\times21$ completely positive matrix $Z$ of rank two with
\[
\|Z\|_{s,q}^2>\|Z\|_{s,s}\|Z\|_{q,q},\qquad (s,q)=(6/5,6).
\]
\end{theorem}
\begin{proof}
Let $X$ have 21 rows, grouped as follows.  For each group use the indicated multiplicity $m_g$, weight $w_g$, and direction $y_g$:
\[
\begin{array}{c|cccc}
g&1&2&3&4\\ \hline
m_g&1&18&1&1\\
w_g&46&17&42&1\\
y_g&(1,0)&(1,1/50)&(1,3/50)&(0,1).
\end{array}
\]
Every row in group $g$ is $w_g^{5/6}y_g$, and set $Z=XX^{\mathsf T}$.  Thus $Z\succeq0$, $Z$ is entrywise nonnegative, and $\operatorname{rank}Z=2$; being a Gram matrix of nonnegative vectors, $Z$ is completely positive.  Put
\[
R_Z=\frac{\|Z\|_{6/5,6}}
{\sqrt{\|Z\|_{6/5,6/5}\|Z\|_{6,6}}}.
\]
Direct interval evaluation gives
\begin{align*}
R_Z&\ge 1.00000061733911154365777590600087863798,\\
R_Z&\le 1.00000061733911154365777590600087863800.
\end{align*}
The entire interval lies above one.  The source package regenerates $Z$ from the four groups and evaluates the ratio independently at 100 digits; a Sage interval script supplies outward rounding.
\end{proof}

\begin{remark}[Rank one is the equality case]\label{rmk:rankone}
Theorem~\ref{thm:mixed} is sharp in the rank.  A symmetric matrix of rank one is $Z=\pm xx^{\mathsf T}$, so $|z_{ij}|=|x_i||x_j|$ and the inner sum factors out of the outer one:
\[
\|Z\|_{s,q}=\Bigl(\sum_j|x_j|^q\Bigl(\sum_i|x_i|^s\Bigr)^{q/s}\Bigr)^{1/q}=\|x\|_s\|x\|_q,
\qquad
\|Z\|_{s,s}=\|x\|_s^2,
\qquad
\|Z\|_{q,q}=\|x\|_q^2 .
\]
Hence $\|Z\|_{s,q}^2=\|Z\|_{s,s}\|Z\|_{q,q}$ for every $1\le s\le q$: at rank one the conjectured inequality is an identity, with no slack anywhere to lose.  The witness of Theorem~\ref{thm:mixed} has rank two, so the failure appears at the first rank that admits any slack at all.
\end{remark}

\begin{remark}[What the two results together settle]
Theorem~\ref{thm:mixed-general} refutes the question as posed, for every non-integer exponent, but its mechanism needs two distinct blocks and so says nothing about $B=A$.  Theorem~\ref{thm:mixed} closes that gap at one point.  Between them the record is: the inequality holds for $p=1$, for every positive integer $p$, for $q=\infty$, and for completely positive $Z$ of order at most three; and it fails at $(6/5,6)$ for order $21$, and at $p=q$ for every non-integer $p$.  The order-three bound is again FitzGerald--Horn: a positive semidefinite matrix of order at most three has $|Z|$ positive semidefinite, and $s\ge n-2$ there.  Details are in \Cref{sec:proofdetails}.
\end{remark}
\end{cxrefutation}


%%% Local Variables:
%%% mode: LaTeX
%%% TeX-master: "../../tex/main"
%%% End:
